% !TEX root = ../main.tex
% !TEX program = pdflatex
\section*{Введение}
\addcontentsline{toc}{section}{Введение}

Быстродействие современного процессора во многом определяется не частотой ядра, а тем, насколько быстро ему удаётся получать данные из памяти. Задержка обращения к оперативной памяти составляет порядка сотни наносекунд, тогда как за это время ядро успевает выполнить сотни операций. Для сокрытия этой задержки между ядром и памятью размещается иерархия кэш-памяти: несколько уровней небольших, но быстрых буферов, хранящих копии недавно использованных данных. Общие рекомендации по оптимизации с учётом кэш-памяти и предвыборки приведены в руководстве Intel~\cite{intel-opt}.

Эффективность программы, в том числе параллельной программы на суперкомпьютере, напрямую зависит от того, насколько её шаблон обращений к памяти согласуется с параметрами кэша: объёмом каждого уровня, размером строки и степенью ассоциативности. Знание этих параметров позволяет правильно выбирать размер блоков при разбиении задач, порядок обхода массивов и выравнивание данных.

Работа выполнялась на ноутбуке MacBook Air с процессором Apple~M2 (4 производительных и 4 энергоэффективных ядра, архитектура ARMv8, размер страницы виртуальной памяти 16~КБ) под управлением macOS. Программа написана на языке C и скомпилирована компилятором Apple clang с уровнем оптимизации \code{-O2}.
